% part: intuitionistic-logic
% chapter: propositions-as-types
% section: types

\documentclass[../../../include/open-logic-section]{subfiles}

\begin{document}

\olfileid{int}{pty}{typ}

\olsection{రకాలు}

$(MN)$ వంటి నిరూపణ పదం ఒక్కటే, దానికి సంబంధించిన నిరూపణను అనన్యంగా నిర్ణయించదు. కారణం: స్వేచ్ఛా చరాలు—అంటే అక్షయాల నిరూపణ పదాలు—సంబంధిత అక్షయాల్లోని \tetoken{సూత్రాలను}{!!{formula}s} తెలపవు. అయితే సందర్భం~$\Gamma$లో వాటిని తెలిపితే, పదం సరైన నిరూపణ పదమైనప్పుడు నిరూపణ అనన్యంగా నిర్ణయించబడుతుంది.

\begin{defn}
నిరూపణ పదం~$N$కు \emph{సందర్భం} అనేది $\Gamma$ జతల సమితి: ప్రతి చరానికి గరిష్ఠంగా ఒక $x : !A$ జతను కేటాయిస్తుంది; $N$లోని ప్రతి స్వేచ్ఛా చరానికీ $\Gamma$లో అలాంటి జత తప్పక ఉంటుంది.
\end{defn}
\sourcecorrection{OLTEPTPTYTYP-005}{మూల నిర్వచనం స్వేచ్ఛా చరాలకే అనన్య రక కేటాయింపును స్పష్టం చేసింది; అదనపు చరాలకూ పరస్పర విరుద్ధ కేటాయింపులు రాకుండా సందర్భాన్ని ఏకవిలువ కేటాయింపుగా స్పష్టం చేశాం. బంధన సందర్భాన్ని విస్తరించే ముందు బద్ధ చరాలకు తాజా పేర్లు ఎంచుకునే ప్రామాణిక పద్ధతినీ పేర్కొన్నాం.}

$\Gamma$లో $x:!A$ జతలు $N$ స్వేచ్ఛా చరాలకు \emph{మాత్రమే} ఉండాలని నిర్వచనం కోరదు. ఉదాహరణకు $\Gamma =
\{x:!A, y:!B, z:!C\}$ అనేది $\pair{x}{y}$కు సందర్భం.

\begin{defn}\ollabel{defn:type}
  $\Gamma$ అనేది~$N$కు సందర్భమని అనుకుందాం. బంధనంతో సందర్భాన్ని విస్తరించే ముందు, అవసరమైతే బద్ధ చరాల పేర్లను మార్చి, సందర్భంలో ఇప్పటికే వాడని తాజా పేర్లు ఎంచుకుంటాం. $\Gamma$కు సంబంధించి~$N$ యొక్క \emph{రకాన్ని} కింది ఆగమన నియమాలతో నిర్వచిస్తాం:
  \begin{enumerate}
  \item $x:!A \in \Gamma$ అయితేనే $x$ రకం $!A$.
  \item $x:!A, \Gamma$కు సంబంధించి $N$ రకం~$!B$ అయితే, $\lambd[\typeof{x}{!A}][N]$ రకం $!A \lif !B$.
  \item $N$ రకం $!A \lif !B$, $M$ రకం~$!A$ అయితే, $(NM)$ రకం~$!B$.
  \item $N$ రకం $!A$, $M$ రకం~$!B$ అయితే, $\pair{N}{M}$ రకం $!A \land !B$.
  \item $N$ రకం~$!A_1 \land !A_2$ అయితే, $\proj{i}{N}$ రకం~$!A_i$.
  \item $N$ రకం $!A$ అయితే, $\inj[!B]{i}{N}$ రకం $i=1$ప్పుడు $!A \lor !B$, $i=2$ప్పుడు $!B \lor
    !A$.
  \item $M$ రకం $!A \lor !B$, $x:!A, \Gamma$కు సంబంధించి $N_1$ రకం $!C$, $y:!B,
  \Gamma$కు సంబంధించి $N_2$ రకం~$!C$ అయితే, $\dcase{M}{x}{N_1}{y}{N_2}$ రకం~$!C$.
  \item $N$ రకం~$\lfalse$ అయితే, $\abort{!A}{N}$ రకం~$!A$.
  \end{enumerate}
\end{defn}
\sourcecorrection{OLTEPTPTYTYP-001}{ఇంజెక్షన్ రక నియమంలో ప్రకటించిన పదం $N$ అయినా మూలం సంబంధం లేని $!M$ను వాడింది; దాన్ని $N$గా సరిచేశాం.}
\sourcecorrection{OLTEPTPTYTYP-002}{వియోగ తొలగింపు ఆధార సందర్భాల్లోని బద్ధ చరాలు $x$, $y$ అయినా మూల పదం $x_1$, $x_2$లను వాడింది; ప్రకటించిన సందర్భాలతో బంధకాలను సమన్వయం చేశాం.}

\begin{prop}
నిరూపణ పదం~$N$కు, $N$ సందర్భం~$\Gamma$కు సంబంధించి రకం ఉంటే, ఆ రకం అనన్యమైనది.
\end{prop}
\sourcecorrection{OLTEPTPTYTYP-003}{ప్రతి ముడి నిరూపణ పదానికి తప్పక ఒక రకం ఉంటుందని మూలం చెప్పింది. పూర్వ విభాగంలోని తప్పుగా రకీకరించిన పదాలే దానికి వ్యతిరేక ఉదాహరణలు. సరైన వాదన రకం ఉన్నప్పుడు దాని అనన్యత; పరిచయంలోని పునర్నిర్మాణ వాదనకూ సరైన పదం అనే షరతును స్పష్టం చేశాం.}

\begin{proof}
  $N$పై ఆగమనంతో.
\end{proof}

$\Gamma$కు సంబంధించి $N$ రకం~$!A$ అయితే $\Gamma \Sequent N : !A$ అని రాస్తాం. \olref{defn:type}లోని నియమాలు పరిచితంగా కనిపించాలి: అవి \Log{tN2i} నియమాలే, ఒక చిన్న తేడాతో—అంతటా సందర్భం~$\Gamma$ ఒకటే. ఈ నిర్వచనాన్ని \Log{tN3i} కలనశాస్త్రంగా వ్యక్తపరచవచ్చు; దాని నియమాలు \olref{tab:tN3ip}లో ఉన్నాయి.
\sourcecorrection{OLTEPTPTYTYP-004}{ఇక్కడ మూలం పూర్వ వ్యవస్థ పట్టికను సూచించింది; ఈ భాగం ఉపయోగించే ఒకే-సందర్భ రక కేటాయింపు పట్టికకు ఇచ్చిన ప్రత్యేక గుర్తుకు సూచనను సమన్వయం చేశాం.}

\subfile{rules-tN3}

నిరూపణ పదాలు $\lambd$-కలనశాస్త్రంలోని ఒక రూపం—\emph{లబ్ధ, యోగ, ఖాళీ రకాలతో కూడిన రక $\lambd$-కలనశాస్త్రం}—లోని పదాలు. సాంప్రదాయ $\lambd$-కలనశాస్త్రంలో చరాలు, \emph{$\lambd$-అమూర్తీకరణలు}~$\lambd[x][N]$, ప్రయోగం $(MN)$తో నిర్మించిన పదాలే ఉంటాయి; రక $\lambd$-అమూర్తీకరణ $\lambd[\typeof{x}{!A}][N]$లోని~$!A$ గుర్తు దానిలో ఉండదు. $\lambd$-కలనశాస్త్రం ట్యూరింగ్-పూర్ణ గణన నమూనా (అంటే ప్రతి పాక్షిక గణనయోగ్య ప్రమేయాన్ని దానిలో ఒక పదంతో ప్రాతినిధ్యం చేయవచ్చు). నిరూపణ పదాలు, రక నియమాలు ఇచ్చే $\lambd$-కలనశాస్త్రం పరిమిత గణన నమూనా; కానీ ప్రాథమిక ఆలోచనలు ఒకటే. రక కేటాయింపు వల్ల పరిమితి వస్తుంది: \emph{సరిగ్గా రకీకరించిన} పదాలే—అంటే సందర్భం~$\Gamma$కు సంబంధించి రకం గల నిరూపణ పదాలే—అనుమతించబడతాయి. ఉదాహరణకు $\lambd[x][(xx)]$ సాంప్రదాయ రకరహిత $\lambd$-కలనశాస్త్రంలో అనుమతించబడుతుంది; కానీ $\lambd[\typeof{x}{!A}][(xx)]$ ఏ $!A$కూ సరిగ్గా రకీకరించబడదు. కారణం: $(xx)$కు రకం ఉండాలంటే, మొదటి~$x$ రకం $!A \lif !B$, రెండవ $x$ రకం~$!A$ కావాలని నియమాలు కోరతాయి. $!A$ అనేది~$!A \lif !B$కు సరైన ఉప-\tetoken{సూత్రం}{!!{formula}} కనుక అది అసాధ్యం.

ఒక రకం పదం ఏ విధమైన వస్తువును సూచిస్తుందో తెలిపేదిగా రకాన్ని సహజంగా భావించవచ్చు. $!A \land !B$ రకం పదం, మొదటి భాగం రకం~$!A$, రెండవ భాగం రకం~$!B$ గల జతను సూచిస్తుంది; అంటే లబ్ధం~$!A \times !B$లో \tetoken{ఒక మూలకం}{!!a{element}}. $!A \lif !B$ రకం పదం, $!A$ రకం వస్తువుల నుంచి~$!B$ రకం వస్తువులకు ప్రమేయం. $!A \lor !B$ రకం పదం $!A$ రకం వస్తువు లేదా~$!B$ రకం వస్తువు, అది ఏదో తెలిపే సమాచారంతో కూడినది. ఇది రెండు సమితులు $!A$, $!B$ల వేరు సమ్మేళనం వంటిది; దాని నిర్వచనం $\{\tuple{i,x}
: i = 1 \text{ లేదా } 2, x \in !A \text{ అయితే } i = 1, x \in !B \text{ అయితే
} i=2\}$. $\lfalse$ రకం ఖాళీ సమితి; అంటే ఆ రకం పదం ఏ వస్తువునూ సూచించలేదు. సహజ నిగమనం (\Log{tN3i}తో సహా) వైరుధ్యరహితమైనది కనుక, తెరిచి ఉన్న ఊహలు లేకుండా~$\lfalse$కు \tetoken{నిరూపణ}{!!{proof}} లేదు; అందువల్ల~$\lfalse$ రకం గల మూసిన నిరూపణ పదం (స్వేచ్ఛా చరాలు లేని పదం) లేదు.

రక $\lambd$-కలనశాస్త్ర సంచాలకాలు, తమకు ఇచ్చిన పదాలు తగిన వస్తువులను సూచిస్తే, తగిన రకం వస్తువులను సూచించే పదాలను ఉత్పత్తి చేస్తాయి. ఉదాహరణకు “$N_1 : !A$, $N_2 : !B$ అయితే $\pair{N_1}{N_2} : !A \land !B$” అంటే: $N_1$ రకం~$!A$ గల వస్తువును, $N_2$ రకం~$!B$ గల వస్తువును సూచిస్తే, $\pair{N_1}{N_2}$ రకం $!A \land !B$ గల వస్తువును సూచిస్తుంది.

\end{document}
